\documentclass[11pt]{article}
\usepackage[T1]{fontenc}
\usepackage{amsmath,amssymb,amsthm,mathrsfs,geometry,hyperref,longtable}
\geometry{margin=26mm}
\setlength{\emergencystretch}{3em}
\newtheorem{theorem}{Theorem}[section]
\newtheorem{lemma}[theorem]{Lemma}
\newtheorem{proposition}[theorem]{Proposition}
\newtheorem{definition}[theorem]{Definition}
\theoremstyle{remark}\newtheorem{remark}[theorem]{Remark}
\newcommand{\ddc}{dd^c}
\newcommand{\PSH}{\operatorname{PSH}}
\newcommand{\cI}{\mathcal I}
\newcommand{\MA}{\operatorname{MA}}
\newcommand{\PP}{\mathbf P^1}
\newcommand{\R}{\mathbf R}
\newcommand{\TV}{\operatorname{TV}}
\title{A geometric polar-measure sector for P14.4\\
\large Reconstruction, modification compatibility, and a birational loss-of-data test}
\author{Unverified research draft}
\date{6 October 2026}
\begin{document}
\maketitle
\begin{abstract}
The unchanged P14.4 construction in arbitrary big non-nef real
Bott--Chern classes is not solved here. We construct an independently
admissible cone of geometric measure data over the valuative carrier
for a separated-current sector on
$\operatorname{Bl}_{(\infty,\infty)}(\PP\times\PP)$.
Its inputs are diffuse positive measures with an explicit logarithmic
kernel condition, not complex potentials. They reconstruct positive
currents, their primitives and ordinary singularity envelopes in this
sector, with order, scaling, constants, actual decreasing approximation
and functorial pullback to every given smooth modification.
The construction retains a diffuse polar current that ideal data and
nonpluripolar products discard. We independently reconstruct the
mandatory equal-ideal/equal-product/equal-positive-mass example and
strengthen its equality to the entire tower of smooth modifications.
This is an analytic geometric sector, not an intrinsic characterization
of a full non-Archimedean cone. No DXZ/BJ/MP operator identification,
arbitrary transcendental extension, or P14.5/P14.6 result is claimed.
\end{abstract}

\section{Scope, conventions, and external inputs}
Let $\ddc=(i/2\pi)\partial\bar\partial$, so that
$\ddc\log|z|^2=\delta_0$. A $\theta$-psh function has the canonical
representative whose sum with each smooth local potential of $\theta$
is psh; upward changes at individual points are excluded. Write
$T_u=\theta+\ddc u$,
$\cI(tu)=\{f:|f|^2e^{-tu}\text{ locally integrable}\}$ for $t>0$,
and $M_\theta(u)=\int_X\langle T_u^2\rangle$.
Masses are raw, not probabilities.
The ordinary envelope is
\[
 P_\theta[u]=\left(\lim_{C\to\infty}P_\theta(u+C,0)\right)^*,
 \quad
 P_\theta(f,g)=
 \left(\sup\{v\in\PSH(X,\theta):v\leq\min(f,g)\}\right)^*.
 \tag{1.1}
\]
For the positive-mass comparison below, ordinary $P$-equivalence
means equality of these normalized envelopes.

Only the complete supplied body B is a caller grant. Its domain is
smooth integral projective varieties over algebraically closed
trivially valued characteristic-zero fields, with actual big rational
bundles. Section~\ref{Bboundary} uses its stationary generated moving
system and raw intersection conventions in that domain only.
The earlier P14.4 paper and catalogue receipts supply no proof grants.

Here are the precise advanced analytic imports used in this paper.
Their complete pertinent statement/proof units were read in the literal
versions below. The ledger identifies local renderer payloads and
their hashes; these hashes are not full-PDF hashes.
\begin{enumerate}
\item BEGZ, arXiv:0812.3674v2, Section 1.2, Definition 1.1,
Propositions 1.4 and 1.6: for local psh $u_j$, the nonpluripolar
product is determined by the increasing truncation currents
\[
 \mathbf1_{\cap_j\{u_j>-N\}}\bigwedge_j
 \ddc\max(u_j,-N).
\]
When defined, it depends only on the currents, is local in the
plurifine topology, and charges no pluripolar set.
It is defined for closed positive $(1,1)$-currents on a compact
K\"ahler manifold. We use locality on ordinary open smooth loci
and the noncharging property. The discussion after Definition 1.2
also gives this local criterion: for potentials locally bounded off
a closed complete pluripolar set, the product is well-defined exactly
when its Bedford--Taylor product on the complement has locally finite
mass; it is then the trivial extension. This criterion applies on our
smooth modifications without assuming they are K\"ahler.
\item DDL, arXiv:1705.05796v3, Theorems 1.1/2.4:
on a compact K\"ahler $n$-fold, for smooth closed real forms
$\theta_j$ with pseudoeffective classes and
$u_j,v_j\in\PSH(X,\theta_j)$, if $v_j\leq u_j+C_j$, then
\[
 \int_X\left\langle\bigwedge_j(\theta_j+\ddc u_j)\right\rangle
 \geq
 \int_X\left\langle\bigwedge_j(\theta_j+\ddc v_j)\right\rangle.
 \tag{1.2}
\]
No small-unbounded-locus hypothesis is required.
\item DDL, Section 3.2, the membership assertion and its proof before
Lemma 3.7: in a big class, if $\phi\leq0$ is $\theta$-psh, then
$\phi\leq P_\theta[\phi]\leq0$ and
$M_\theta(P_\theta[\phi])=M_\theta(\phi)$.
We use this assertion only when $M_\theta(\phi)>0$.
For clarity, the stronger Theorem 3.12, whose complete proof was also
read, says in that positive-mass range that the envelope is the
maximum of
$\{v:\phi\leq v\leq0,\ M_\theta(v)=M_\theta(\phi)\}$ and is idempotent.
We do not need that stronger characterization.
\end{enumerate}
The source DDL uses $i\partial\bar\partial$; applying its assertions to
$u/(2\pi)$ and scaling all obstacles and constants by $1/(2\pi)$
gives exactly the stated convention, with the same currents.
Its Section 3.2 membership proof uses Theorem 2.3 (capacity lower
semicontinuity, and weak convergence when limiting mass dominates
the limsup) and Theorem 2.4; their complete proof units were read.
We invoke the published theorems, not a purported independent proof
of their cited background. No source theorem concerning DXZ, MP,
analytic Zariski approximation, or projective resolution is imported.

\section{Independent geometric data on a curve}
Fix the unit Fubini--Study form $\omega$ on $\PP$ and its affine
coordinate $z$. For finite $t$ set
\[
 G(z,t)=\log\frac{|z-t|^2}{(1+|z|^2)(1+|t|^2)}\leq0,\qquad
 \ddc_zG(z,t)=\delta_t-\omega.
 \tag{2.1}
\]
The formula extends at infinity in coordinate $1/z$. Its sign follows
from $|z-t|^2\leq(1+|z|^2)(1+|t|^2)$.

\begin{definition}
In this paper an admissible measure is a finite positive Borel measure
$\nu$ on $\PP$, with no atoms, whose support $S$ is a compact subset
of the finite affine chart, and which satisfies
\[
 \int_\PP G(x,t)\,d\nu(t)=-\infty\quad\text{for every }x\in S.
 \tag{2.2}
\]
The zero measure is admissible, with empty support.
\end{definition}
This is a condition on a measure and a fixed geometric kernel.
It presupposes neither a complex psh input nor a $P$-envelope.
It is restrictive: arbitrary polar measures or arbitrary singularities
are not asserted to satisfy it. Positive sums and positive scalar
multiples are admissible. Indeed their supports are the finite union
of the incoming supports, and the nonpositive kernel makes an infinite
negative contribution at every point of that union.

\begin{lemma}\label{curvelemma}
For an admissible $\nu$ of mass $m$, the integral
$U_\nu(z)=\int G(z,t)d\nu(t)$ is nonpositive, locally integrable,
has the canonical quasi-psh representative, and
$\ddc U_\nu=\nu-m\omega$. It is smooth and finite off $S$,
and its pole set is precisely $S$. For every $h>0$,
$e^{-hU_\nu}$ is locally integrable.
\end{lemma}
\begin{proof}
The local $L^1$ norms of logarithmic kernels with $t$ in a compact
set are uniformly bounded, by translating their singularities into
a fixed larger disc. Fubini proves integrability and the Hessian
identity. Fatou applied to $-G$ proves upper semicontinuity.
Integrating the submean inequalities for $\log|z-t|^2$ proves the
canonical local subharmonic condition after adding $m\log(1+|z|^2)$.
Away from the support all kernel derivatives are locally uniformly
bounded. Condition (2.2) identifies the poles. The function is finite
near infinity, hence nontrivial. In particular $S$ is polar and has
area zero, since it is contained in the infinite locus of a locally
integrable function.

At a point $z_0\in S$, atomlessness lets us choose a disc $B(z_0,r)$
with mass $\delta<1/h$. On the smaller disc $B(z_0,r/2)$ the
complementary potential and spherical terms are bounded.
If $\delta>0$, Jensen for the normalized restricted measure gives
\[
 \exp\left(-h\int_{B(z_0,r)}\log|z-t|^2\,d\nu(t)\right)
 \leq\frac1\delta\int_{B(z_0,r)}|z-t|^{-2h\delta}\,d\nu(t).
 \tag{2.3}
\]
Truncating the logarithm first justifies the extended-value formula.
The integrals of the right side on the smaller disc are uniformly
bounded by $\int_{|w|<2r}|w|^{-2h\delta}dA(w)<\infty$.
Fubini proves the assertion. The case $\delta=0$ and points off
the support are immediate.
\end{proof}

\begin{lemma}[Pullback integrability without a birational grant]\label{pulllemma}
If $f$ is a nonconstant holomorphic map from a connected complex
manifold to $\PP$, then $U_\nu\circ f$ has its canonical quasi-psh
pullback representative and $e^{-hU_\nu\circ f}$ is locally
integrable for every $h>0$.
\end{lemma}
\begin{proof}
Near a point with image outside $S$, including infinity, this is
smooth. At a point with image $t_0\in S$, choose a local complex
line direction in which $f(z,0)-t_0$ has a finite order $d\geq1$.
Such a direction exists by the first nonzero homogeneous term of
the holomorphic Taylor series. Choose a small circle $|z|=R$
on which this function has no zero, with all its $d$ zeros
strictly inside $R/2$. For nearby other coordinates $w$ and $t$,
the argument principle, or its homotopy proof, gives exactly $d$
zeros $r_i(w,t)$ inside $R/2$.
Factoring them, including multiplicities, gives
\[
 f(z,w)-t=u(z,w,t)\prod_{i=1}^d(z-r_i(w,t)).
 \tag{2.4}
\]
For each fixed $w,t$, $u$ is holomorphic and zero-free on the disc.
On the boundary its modulus is uniformly bounded below: the numerator
is bounded below and the product is bounded above.
The maximum principle for $1/u$ gives the same lower bound inside.
This argument does not require holomorphic choices of the roots.
H\"older therefore gives a uniform integral bound on a smaller disc
for $|f(z,w)-t|^{-2\epsilon}$ whenever $d\epsilon<1$, using
the $d$ factors and the bounds for
$\int|z-r|^{-2d\epsilon}dA(z)$.
Integrating over $w$ gives the corresponding local bound.

Choose the restricted measure near $t_0$ to have
$\delta<1/(hd)$, possible by atomlessness. Jensen as in (2.3),
now with $z$ replaced by $f(z,w)$, and the uniform bound just proved
give the required exponential integrability.
The complementary part is bounded after shrinking the image.
The same factorization bounds local logarithmic $L^1$ norms
uniformly in $t$, so Fubini defines the pullback primitive.
Alternatively its psh local primitives are compositions of psh
functions with $f$, verified by their restriction to holomorphic
discs and the submean inequality. They are nontrivial: a nonconstant
holomorphic scalar has open local image, whereas $S$ has no interior.
These arguments also establish the canonical representative.
\end{proof}

\section{An enriched carrier and a geometric sector cone}
Let $Y=\PP_z\times\PP_w$ and
$\pi:X=\operatorname{Bl}_{p}Y\to Y$, $p=(\infty,\infty)$,
with exceptional curve $E$. Put
$q=\operatorname{pr}_z\circ\pi$, $r=\operatorname{pr}_w\circ\pi$.
Let $h_E$ be a smooth metric, $\rho=c_1(\mathcal O(E),h_E)$,
$e=\log|\sigma_E|_{h_E}^2\leq0$, so $\ddc e=[E]-\rho$.
Define $H_z=\{q^*\omega_z\}$ and $H_w=\{r^*\omega_w\}$.

Let $\mathscr V_X$ be the centered coherent-ideal semivaluation
carrier appearing in the question, retaining supports and
$+\infty$. We use only this specified point set and evaluations
$v(\mathcal I_E)$, not a theorem identifying a topology or a
DXZ image. Let $v_0$ be the generic trivial semivaluation.
The precise enriched carrier used here is
\[
 \mathscr C_X=\mathscr V_X\times\PP(\mathbf C),\qquad
 \operatorname{pr}_1:\mathscr C_X\longrightarrow\mathscr V_X.
 \tag{3.1}
\]
The additional geometric datum is the positive measure $\nu$ on the
slice $\{v_0\}\times\PP$, or equivalently its positive functional
$F_\nu(g)=\int g\,d\nu$ on $C(\PP)$.
Using all continuous tests determines a finite Borel measure:
for a compact metric space, continuous cutoffs approximate indicators
of closed sets, and then regularity determines Borel sets.
No potential on $X$ is part of the input.
This deliberately adds complex geometry; it is not a claim that the
MP spectrum alone contains that geometry.

Define $\mathscr D$ by the independently checkable inputs
\[
 D=(a,b,c,\nu,\kappa),\quad a,b,c\geq0,\quad
 m:=\nu(\PP)\leq b,\quad \nu\text{ admissible},\quad\kappa\in\R.
 \tag{3.2}
\]
Addition adds all five entries, positive scaling scales all entries,
and the zero element has $\nu=0$.
For the big non-nef sector we require $a,b,c>0$.
The geometric output, not an input, is
\[
 \begin{split}
 \alpha_D&=bH_z+cH_w+a\{E\},\\
 \theta_D&=bq^*\omega_z+cr^*\omega_w+a\rho,\\
 \Phi_D&=ae+U_\nu\circ q+\kappa,\\
 T_D&=a[E]+(b-m)q^*\omega_z+cr^*\omega_w+q^*\nu.
 \end{split}\tag{3.3}
\]
Here $q^*\nu$ is defined by the local psh primitives of
Lemma~\ref{pulllemma}: $q^*\nu=\ddc(U_\nu\circ q)+m q^*\omega_z$.
It is positive, not an unjustified pullback of a distribution.

\begin{theorem}[Sector realization]\label{sectorthm}
The data (3.2) form a convex geometric cone with a constants action.
For $a,b,c>0$, $\alpha_D$ is big and non-nef and
$\Phi_D\in\PSH(X,\theta_D)$.
On a fixed $(a,b,c)$ leaf, the exact order is the kernel-test order
\[
 D\leq D'\quad\Longleftrightarrow\quad
 U_\nu(z)+\kappa\leq U_{\nu'}(z)+\kappa'
 \quad\text{for every }z\in\PP.
 \tag{3.4}
\]
It is equivalent to $\Phi_D\leq\Phi_{D'}$ on $X$.
The state determines its ordinary $P$-envelope by (1.1);
this does not characterize the image of that envelope operator.
Addition and positive scaling of data give the corresponding sums
and scalings of actual currents and primitives.
The measure is reconstructible from $T_D$ and the fixed geometric
parameters, independently of $\kappa$.
\end{theorem}
\begin{proof}
The cone assertions follow from the sign of $G$ and support unions.
Lemma~\ref{curvelemma} and the local psh pullback prove (3.3).
All coefficients of the displayed positive current are nonnegative.

For the geometry, let $H=H_z+H_w$.
The class $A=2H-\{E\}$ is ample: the Segre embedding
$Y\hookrightarrow\mathbf P^3$ and the three hyperplanes through $p$
embed its blowup as the graph in $Y\times\mathbf P^2$.
The hyperplane pullbacks are $H$ and $H-\{E\}$, and their
Segre sum is $A$. Locally the three hyperplanes generate the
ideal of $p$, since their differentials generate its cotangent space;
outside $p$ they have no common zero. This verifies the graph
description by the local Rees algebra of that ideal.
For $t=\min(b,c)>0$,
\[
 \alpha_D=\frac t2 A+(b-t)H_z+(c-t)H_w+
                 (a+t/2)\{E\}.
 \tag{3.5}
\]
A K\"ahler representative of the first term, the semipositive
pullback forms, and the effective exceptional current give a
K\"ahler current in $\alpha_D$.
The exceptional normal bundle is $\mathcal O_{\PP}(-1)$, hence
$E^2=-1$, while $H_zE=H_wE=0$; thus $\alpha_DE=-a<0$.
Nef classes have nonnegative degrees on curves by their
K\"ahler approximants, so this class is non-nef.

Off $E$ the common term $ae$ is finite, and $q$ has points off $E$
above every point of $\PP$. This proves both directions in (3.4);
on $E$ both primitives are $-\infty$ when $a>0$.
If $a=0$, the term $ae$ is defined as zero, not $0(-\infty)$.
Linearity of kernel integrals proves the operation assertions
at finite points and with literal negative infinities at poles.
For zero scaling the primitive is zero, by convention.

On the product away from $p$, integrate $T_D-a[E]$ against
the unit form in the $w$ direction. The resulting measure in $z$
is $(b-m)\omega_z+\nu$.
Its singular part relative to area is exactly $\nu$, since the
support of $\nu$ has area zero. The computation is legitimate
on neighborhoods of $S\times\PP$, which avoid $p$.
Thus the geometric current recovers $\nu$.
In particular, on a fixed leaf different measures cannot give
the same current. The constants coordinate fixes the explicitly
reconstructed primitive, not merely its current.
\end{proof}

\section{All given smooth modifications and reference covariance}
\begin{theorem}\label{modthm}
For every given smooth connected proper bimeromorphic
$p:W\to X$ and every $D\in\mathscr D$ with $a,b,c>0$,
the actual pullback primitive is $\Phi_{D,W}=\Phi_D\circ p$ and
\[
 T_{D,W}=a[p^*E]+(b-m)(q\circ p)^*\omega_z+
 c(r\circ p)^*\omega_w+(q\circ p)^*\nu.
 \tag{4.1}
\]
These data are compatible with further smooth modifications.
No existence or projectivity of a resolution is needed for this
assertion: it quantifies over each modification already given.
\end{theorem}
\begin{proof}
The total pullback $p^*E$ is the zero divisor of the actual pulled
section. In local frames, $\ddc\log|p^*\sigma_E|^2=[p^*E]$
with its full multiplicities. Lemma~\ref{pulllemma} applies to
$q\circ p$, a nonconstant holomorphic map. Adding the pulled
smooth local potentials proves the canonical psh condition and
(4.1). For a second modification composition of functions,
pulled sections and smooth forms is literal.
The definition of the polar current by those same local primitives
then proves composition compatibility. There is no subtraction of
two infinite values and no replacement of $p^*E$ by its strict transform.
\end{proof}

At the point-set carrier level the compatible map is
$\mathscr C_W\to\mathscr C_X$, $(v,z)\mapsto(p_*v,z)$, where
$(p_*v)(\mathfrak a)=v(\mathfrak a\mathcal O_W)$.
Pullback preserves ideal sums, products and units, so this is a
centered semivaluation with its support retained by contraction.
It composes under further modifications, sends the generic trivial
point to the generic trivial point, and keeps the geometric measure
on that slice unchanged. Moreover
$v(\mathcal I_{p^*E})=(p_*v)(\mathcal I_E)$.
This is not a claim about a new identification of MP topologies.

The data are actual geometric data, so changing a smooth representative
by $\theta'=\theta+\ddc f$, $f$ smooth, transports the primitive
to $\Phi'=\Phi-f$ and keeps $T$ unchanged.
Changing the exceptional metric is a special case.
All multiplier ideals remain unchanged by this bounded perturbation.
For a transported reference obstacle $h-f$, rooftop competitors
give the exact formula
\[
 P_{\theta'}(\Phi-f+C,h-f)=P_\theta(\Phi+C,h)-f.
 \tag{4.2}
\]
Limits and upper semicontinuous regularization commute with the
smooth subtraction. Thus the singularity-envelope construction is
covariant with simultaneous transport of the reference.
Resetting the obstacle to zero after changing $\theta$ is a different
reference choice, not the equality in (4.2).
The cone was defined using fixed geometric forms and a fibration;
we do not assert independence of those actual geometric data,
or perturbation independence of a full real-class NA cone.
On this fixed sector no auxiliary ample perturbation enters (3.2)--(4.1).

\section{Projection, products, mass, and an actual approximation}
\begin{theorem}\label{shadowthm}
For $D$ in the big non-nef sector and $t>0$,
\[
 \cI(t\Phi_D)=\mathcal O_X(-\lfloor ta\rfloor E),\qquad
 Q_D(v):=-\lim_{\ell\to\infty}
   \ell^{-1}v(\cI(\ell\Phi_D))=-a\,v(\mathcal I_E).
 \tag{5.1}
\]
The formula includes supported semivaluations. The normalized
ideal shadow has $Q_D(v_0)=0$ and forgets $\nu$ and $\kappa$.
For any smooth modification in Theorem~\ref{modthm},
\[
 \cI(t\Phi_{D,W})=\cI(ta\,e\circ p).
 \tag{5.2}
\]
If $p^*E$ is simple normal crossing, the latter is
$\mathcal O_W(-\lfloor ta\,p^*E\rfloor)$, componentwise.
No such formula is asserted for a non-normal-crossing divisor.

With $S_D=(b-m)q^*\omega_z+cr^*\omega_w$, for $k=1,2$,
\[
 \langle T_D^k\rangle=S_D^k,\quad
 \langle T_{D,W}^k\rangle=(p^*S_D)^k,\quad
 M_{\theta_D}(\Phi_D)=2c(b-m).
 \tag{5.3}
\]
The modification has that same top raw mass.
\end{theorem}
\begin{proof}
Near $E$, $q$ is near infinity and $U_\nu\circ q$ is smooth.
For a defining coordinate $\zeta$, $\Phi_D=a\log|\zeta|^2+O(1)$.
A monomial $\zeta^j$ is integrable exactly when $j+1>ta$.
The least integer is $\lfloor ta\rfloor$, including integral $ta$.
For arbitrary holomorphic germs expand in $\zeta$: a first nonzero
coefficient below that degree is bounded below on an open subset
of the other coordinate, where the radial integral diverges.
Away from $E$, exponential integrability in
Lemma~\ref{curvelemma}, and Fubini, give the unit ideal.
This proves the sheaf equality.
For $w=v(\mathcal I_E)<\infty$, the ideal value is
$\lfloor\ell a\rfloor w$. For $w=\infty$, it is infinity for
all sufficiently large $\ell$. The unit ideal has value zero,
so no $0\cdot\infty$ is evaluated. This proves (5.1) on every
specified semivaluation, not by divisorial density.

On $W$, $p^{-1}E$ maps under $q\circ p$ to infinity.
The polar term is therefore bounded near that divisor.
Off it, Lemma~\ref{pulllemma} makes all exponentials of the
polar term locally integrable, and the exceptional term is smooth.
This proves (5.2) locally everywhere. In an SNC chart the
logarithmic divisor is a sum of coordinate logarithms, and
separate radial integrals give the componentwise floor formula.

The closed pole set on $W$ is
$Z_W=p^{-1}E\cup(q\circ p)^{-1}S$.
It is pluripolar by the actual quasi-psh primitive (when $\nu=0$
only the first set is required).
Off it the current is the smooth form $p^*S_D$.
BEGZ locality and noncharging give (5.3) if the smooth forms
have no coefficient mass on $Z_W$.
For the analytic divisor this is immediate from area zero.
For $(q\circ p)^{-1}S$, outside the critical locus use a
holomorphic submersion coordinate and Fubini; $S$ has area zero.
The critical locus has area zero: some derivative of the
nonconstant holomorphic scalar is not identically zero, so this
locus is contained in the zero set of a nonzero holomorphic function.
Such a zero set has area zero by slicing and the one-variable
isolated-zero theorem. This proves the assertion also at critical points.

The forms from either curve have square zero, so
$S_D^2=2c(b-m)\pi^*(\omega_z\wedge\omega_w)$.
The last volume form has integral one, by the isomorphism off
$E$ and $p$. A bimeromorphic map between smooth surfaces is
an isomorphism off proper analytic subsets; changing variables
there shows $\int_W(p^*S_D)^2=\int_XS_D^2$.
This proves both raw mass assertions. BEGZ applies on $W$:
one may use its local definition and finite mass directly here,
so no theorem that an arbitrary smooth modification is K\"ahler
is needed. The smooth extension calculation already proves
well-definedness and locality suffices.
\end{proof}

For fixed $a,b,c$ the ideal-shadow projection is constant, so it
preserves order trivially and cannot reflect it. Positive scaling
and addition descend to its asymptotic profiles $Q$, with the
literal extended-value conventions above. Exact multiplier ideals
do not have an additive floor law:
$\lfloor t(a+a')\rfloor$ need not be
$\lfloor ta\rfloor+\lfloor ta'\rfloor$.
Constants act trivially on ideals; the separate coordinate
$\kappa$ acts on the primitive by addition. There is no assertion
that $\kappa$ is the endpoint of a compatible DXZ test curve.
Neither the raw mass (5.3), the polar measure, nor the ordinary
$P$-envelope descends to the ideal shadow.

\begin{proposition}[Replacement topology and decreasing approximation]\label{approx}
For every fixed $D$, set
\[
 \nu_j=(1-1/j)\nu,\quad
 D_j=(a,b,c,\nu_j,\kappa),\qquad j\geq2.
 \tag{5.4}
\]
These are actual admissible states in the same class,
$\Phi_{D_j}\downarrow\Phi_D$ everywhere, and their raw products
and masses converge to (5.3). More generally, on states with
support in a common compact affine set, total variation convergence
of measures implies local $L^1$ convergence of their reconstructed
primitives on $X$ and every fixed smooth modification.
This is a topology on the admissible sector, not a formal completion.
\end{proposition}
\begin{proof}
$G\leq0$ gives the pointwise decrease, including at all poles.
The measure error is $m/j$ in total variation. The smooth products
depend polynomially on $b-m_j,c$, where $m_j=(1-1/j)m$;
this proves product and mass convergence without a general
Monge--Amp\`ere continuity assertion.
Uniform local $L^1$ kernel bounds give
$\|U_\nu-U_{\nu'}\|_{L^1(K)}\leq C_K\|\nu-\nu'\|_{\TV}$
on the curve. On a fixed modification the uniform logarithmic
bounds from the proof of Lemma~\ref{pulllemma}, and a finite chart
cover, give the same estimate for the composed kernels.
Constants and the smooth parameters can be transported separately.
No density of finite ideals or ordinary PL models is asserted.
Admissibility is not asserted to be closed in arbitrary weak
measure limits.
\end{proof}

At positive $\lambda$, competitor scaling gives
$P_{\lambda\theta}[\lambda u]=\lambda P_\theta[u]$;
constant shifts of $u$ leave $P_\theta[u]$ unchanged by reindexing
$C$. Rooftop comparison proves monotonicity of this envelope
map. The current and raw mass scale by $\lambda$ and $\lambda^2$,
respectively. These are sector statements about actual complex
currents and ordinary envelopes.

\section{Explicit measures and separation on the entire birational tower}
Set $\ell_0=1$, $\ell_j=e^{-4^j}$ for $j\geq1$.
At every level retain the two closed intervals of the next length
at the endpoints of every retained interval, starting from $[0,1]$.
The inequality $2\ell_{j+1}<\ell_j$ makes the children disjoint.
The intersection is a compact set $C$.
The Bernoulli measure $\mu_0$ assigns each level-$j$ interval
mass $2^{-j}$; its singletons have mass at most $2^{-j}$ for
every $j$, so it is atomless. Let $D=C+2$ and $\mu_1$ be the
translated probability.
For $x$ in its level-$j$ interval,
\[
 U_{\mu_0}(x)\leq 2^{-j}\,2\log\ell_j=-2\,2^j\longrightarrow-\infty.
 \tag{6.1}
\]
This uses $G\leq0$ on the complement and denominators at least one.
The same proof holds on $D$. Thus both measures are admissible,
their supports are disjoint, and all curve and pullback
integrability assertions above apply.

\begin{theorem}[Mandatory test, strengthened]\label{mandatory}
Fix $a,b,c>0$ and $0<s<b$. In the fixed class
$\alpha=bH_z+cH_w+a\{E\}$ use the states
$D_i=(a,b,c,s\mu_i,0)$, $i=0,1$.
Their complete real-scale multiplier-ideal families coincide on
$X$ and on every given smooth modification $W\to X$.
Their entire first and second nonpluripolar currents coincide
on every such $W$, and their positive raw mass is
$2c(b-s)$.
Nevertheless their normalized ordinary $P$-envelopes on $X$ differ.
The enriched measure data distinguish them.
\end{theorem}
\begin{proof}
The equalities are (5.1)--(5.3), since the parameters and masses
of the two measures agree. The measures themselves differ,
as they have disjoint supports and positive mass.

The maximum of $U_{\mu_0}$ and $U_{\mu_1}$ is bounded:
near either compact support the other function is smooth and
bounded below, and on the remaining compact set both are smooth.
Thus
$W_0=\max(\Phi_{D_0},\Phi_{D_1})$
differs by a bounded function from $ae$.
It is $\theta$-psh by the local psh maximum rule.
Mass monotonicity (1.2) in both directions, and the $\nu=0$
case of (5.3), give $M_\theta(W_0)=2bc$.
Both inputs are nonpositive and have strictly positive mass.
If their ordinary envelopes were a common $P$, the DDL membership
assertion would give $\Phi_{D_i}\leq P$ and
$M_\theta(P)=2c(b-s)$.
Then $W_0\leq P$ and (1.2) gives
$2bc\leq2c(b-s)<2bc$, a contradiction.
Every DDL hypothesis is now checked: $X$ is compact K\"ahler,
the class is big, all competitors are $\theta$-psh,
and the two inputs have positive mass. No zero-mass envelope
assertion and no commutation of envelopes with maxima was used.
\end{proof}

In particular, $b=c=1$, $a=s=1/2$ gives the precise mandatory
class $\alpha=H+\frac12E$ with mass one,
$\cI(t\Phi_i)=\mathcal O_X(-\lfloor t/2\rfloor E)$,
first product $\frac12q^*\omega_z+r^*\omega_w$, and top measure
$\pi^*(\omega_z\wedge\omega_w)$.
The ordinary self-intersection is $2-a^2=7/4$, not the raw mass.
The strengthened conclusion concerns only the specified birational
ideal/product tower. It does not assert equality of every conceivable
movable datum, and is not a negative answer to original P14.4.
The full polar currents $(q\circ p)^*\nu$, retained here,
already escape this equality.

\section{The actual moving system permitted by B}\label{Bboundary}
For rational $a,b,c>0$, let
$L=\pi^*\mathcal O_Y(b,c)+aE$ on the complex projective surface,
viewed over the trivially valued field $\mathbf C$.
For cleared degrees $\ell$, every section of $\ell L$ restricts
to a line bundle of negative degree $-\ell a$ on $E$, so vanishes
there. Repeating this argument after each division gives
\[
 H^0(X,\ell L)=\sigma_E^{\ell a}
      H^0(X,\pi^*\mathcal O_Y(\ell b,\ell c)).
 \tag{7.1}
\]
The remaining space descends to $Y$.
One elementary justification is to view it as rational sections:
a pole downstairs along a prime divisor would remain a pole
upstairs, and a holomorphic section on $Y\setminus\{p\}$ extends
across the smooth surface point by the local Hartogs expansion.
The remaining bundle is generated since both degrees are positive.
Thus the actual full moving pair is
$A=\pi^*\mathcal O_Y(b,c)$, $T=aE$.
This is an actual positive part, not a relatively nef lift of $L$.

In B's permitted rational-bundle domain the full zero-weight
section system has the genuine all-point scalar
$q_T(v)=-a\,v(\mathcal I_E)$.
At each centered analytic semivaluation some moving generator is
a unit, since the system is generated; the equality includes
supported points. Every normalized FS scalar on the fixed $L$
is at most $q_T$, by (7.1) and its nonpositive weights.
For B's varying-bundle endpoints the same inequality holds:
at a fixed divisorial valuation the degree computation forces
a vanishing coefficient at least
$\delta_i=-c_1(B_i)\cdot E$ for a nearby actual bundle $B_i$,
up to degree-rounding error. After removing $kE$, the restriction
degree is $\ell B_i\cdot E+k$, and it remains negative until
$k\geq-\ell B_i\cdot E$. Thus $\delta_i\to a$.
The remaining canonical section norms are at most one, and the
incoming generic constants tend to zero. At each fixed divisorial
$v$ this bounds the endpoint by $-a v(E)$.
B's all-point order determination, with its $v_0$ coordinate,
gives the inequality everywhere.
Since the zero-weight system itself realizes $q_T$,
it is B's normalized minimal reference in this sector.

For completeness the product moving model $X\times\mathbf P^1$
with $C=\operatorname{pr}_X^*A+dF$, $d>0$ rational, has
\[
 A^2=2bc,\qquad C^3=3d(2bc),\qquad
 e_A(0)=C^3/3-dA^2=0.
 \tag{7.2}
\]
Its sole central component has multiplicity one, valuation $v_0$,
and restriction square $2bc$. B's exact finite component formula
therefore gives the raw measure $2bc\,\delta_{v_0}$.
The physical horizontal factor is still $aE$.
This is a calculation of the permitted operator, not only a
numerical volume identity. Positive scaling by rational bundle
powers gives the displayed raw normalization.
The complex state with $\nu=0$ has raw measure
$2bc\,\pi^*(\omega_z\wedge\omega_w)$ on a different carrier.
We assert no equality of these two measures or canonical
pushforward between those carriers.

For $\nu=s\mu_i$ with $s>0$, the identical ideal profile still
equals $q_T$, but its complex raw mass is $2c(b-s)$.
Consequently B cannot supply a $P$-sensitive complex comparison
by identifying these single-weight ideal profiles with its reference
operator. These are not compatible DXZ test curves, and no
non-Archimedean operator is defined by transporting that complex mass.
Grant B gives neither this transport nor an extension to real
transcendental or nonprojective classes. Its internal G0/G1 are
not invoked as separate grants.

\section{Precise unsolved requirements and evidence}
This construction is more than appending an original complex potential:
its admissibility is (2.2), its cone is (3.2), its geometric state is
a positive measure functional, and reconstruction and modification
compatibility have been proved. It is nonetheless only a sector
construction with prescribed fibration and geometric forms.
Its real coefficients lie in the real N\'eron--Severi space of
this projective surface; irrational coefficients do not establish
a genuinely transcendental case.
It is not the independently characterized full NA positivity theory
asked for in unchanged P14.4.

There is even a concrete missing cone operation in this sector:
the maximum of the two reconstructed curve weights in
Theorem~\ref{mandatory} is bounded and nonconstant.
It cannot be $U_\nu+\kappa$ for an admissible nonzero $\nu$,
which has poles, or for zero $\nu$, which is constant.
To see nonconstancy, near a support point one weight is very negative
while the other is smooth; on a nonempty ordinary open set the
winner has Hessian $-\omega$, hence cannot be constant.
Thus the sector is not closed under maxima in a fixed class.
Its decreasing approximation is (5.4), not a decreasing PL or
moving-model approximation theorem.

The following requirements remain unresolved here:
\begin{enumerate}
\item An intrinsic admissible geometric cone for every compact
K\"ahler manifold and every big non-nef real Bott--Chern class,
including genuinely transcendental and nonprojective classes.
\item Reconstruction of arbitrary prescribed ordinary $P$-types,
independent of fibration and kernel choices, with the requisite
maxima/envelopes and full approximation theory.
\item A projection from such states to compatible DXZ objects or
the actual MP potential cone, with all-point comparison and the
required endpoint and perturbation conventions.
\item Positive-volume operator comparison with DXZ/BJ/MP in
genuine common domains, retaining raw singularity-dependent mass
and the correct reference. No energy extension is attempted here.
\end{enumerate}
Known full K\"ahler image, upper semicontinuity, and strict-interior
filtration maximality from the catalogue are not relisted as open;
we simply do not use them as unsupplied proofs. No other active
P14 atom is claimed solved.
The relative-nef generic-fibre obstruction is not avoided by a
false nef total-space lift: (7.1) explicitly keeps its horizontal
fixed factor and generated moving part. No analogue of that full
section-system decomposition in arbitrary real analytic classes
is established here.

All artifacts are local. The complete parallel argument is
\texttt{P14.4-argument.md}; structured notes are
\texttt{supporting-notes.json}; the quote-free application ledger is
\texttt{source-application-\allowbreak ledger.json}. Literal renderer ranges
cover all used advanced statement/proof units without gaps.
Direct PDF retrieval failed at the supplied proxy; web-rendered
literal-version payloads were used without changing any configuration.
Their byte hashes certify local integrity only.
The protected external editor was not accessed.
No additional model route, delegated agent, installation or
resource-heavy computation was used.
Rerun \texttt{python check\_evidence.py}, then
\texttt{latexmk -pdf -interaction=nonstopmode -halt-on-error
P14.4-paper.tex}. Compile success and finite arithmetic diagnostics
are not mathematical certification. The manuscript and proofs
remain unverified drafts; the original atom remains open.

\begin{thebibliography}{9}
\bibitem{BEGZ} S. Boucksom, P. Eyssidieux, V. Guedj and A. Zeriahi,
\emph{Monge--Amp\`ere equations in big cohomology classes},
Acta Math. 205 (2010), 199--262.
Literal version \href{https://arxiv.org/pdf/0812.3674v2}{arXiv:0812.3674v2}.
\bibitem{DDL} T. Darvas, E. Di Nezza and C. H. Lu,
\emph{Monotonicity of non-pluripolar products and complex
Monge--Amp\`ere equations with prescribed singularity},
Analysis \& PDE 11 (2018), 2049--2087.
Literal version \href{https://arxiv.org/pdf/1705.05796v3}{arXiv:1705.05796v3}.
\end{thebibliography}
\end{document}
